\documentclass[10pt,a4paper]{amsart}
\usepackage[margin=19mm]{geometry}
\usepackage{amsmath,amssymb,mathtools}
\usepackage[scr=rsfs,cal=euler]{mathalfa}
\usepackage{xfrac}
\usepackage[unicode,hidelinks,bookmarks=false]{hyperref}
\newcommand{\ie}{i.e.\ }
\newcommand{\cf}{cf.\ }
\newcommand{\resp}{respectively\ }
\newcommand{\Subsection}{\subsection}
\providecommand{\nobreakdash}{\nobreak\hbox{-}}
\setlength{\emergencystretch}{2em}
\multlinegap0pt
\usepackage{bm}
\usepackage{bbm}
\usepackage{dsfont}
\usepackage{xspace}
\usepackage{cancel}
\usepackage{mathtools}
\usepackage{amsthm}
\makeatletter
\@ifundefined{Th}{\newtheorem{Th}{Th}}{}
\@ifundefined{Prop}{\newtheorem{Prop}[Th]{Proposition}}{}
\makeatother
\makeatletter
\expandafter\ifx\csname Proof\endcsname\relax
\newenvironment{Proof}
{\par\noindent{\bf Proof.}}


\fi
\newcommand{\ex}{\mathbb{E}}
\makeatother
\title[Drift Control with Discretionary Stopping for a Diffusion]{Drift Control with Discretionary Stopping for a Diffusion}
\author{V\'aclav E. Bene\v{s}, Georgy Gaitsgori, Ioannis Karatzas}
\date{Source: arXiv:2401.10043v3 (CC BY 4.0)}
\begin{document}
\maketitle

\noindent\emph{Source and licence.} Drift Control with Discretionary Stopping for a Diffusion. Version arXiv:2401.10043v3 (CC BY 4.0), distributed under the Creative Commons Attribution 4.0 International licence, as recorded in the arXiv licence metadata for this exact version and shown on the abstract page https://arxiv.org/abs/2401.10043v3. Full rights notice: licenses/arxiv-2401.10043-CC-BY-4.0.txt.\par\smallskip

\noindent\textit{Editorial scope.} Counted unit: the Prop whose source label is \texttt{prop\_duality\_gap} in the article (source0.tex lines 625--636, source label \texttt{prop\_duality\_gap}), delivered together with the article's own complete original proof of it (source0.tex lines 638--681, one \texttt{proof} environment of 44 lines). No mathematics is written, repaired or re-derived: statement and proof are the article's own text line for line. Required context retained uncounted: value\_function (lines 199--201); value\_function\_constrained (lines 541--543); duality\_gap (lines 572--594). Every definition, display and label of those lines is retained unchanged. Omitted from this package and not redistributed: 1--198, 202--540, 544--571, 595--624, 637--637, 682--979. Nothing inside a retained excerpt is omitted or altered: every retained body is delivered exactly as the article's lines are written, so no omission receipt of any kind accompanies this package. Citations: the retained excerpts carry no citation command at all, so no bibliography entry of the article is transcribed and no reference is redistributed. Reference adaptation: none was needed and none was made; every reference inside the delivered unit resolves inside it, so no unresolved reference or citation is produced. Printed slips kept literal: no printed word, symbol, hypothesis or number of the counted statement or of its proof was altered. Layout and class: the article's own class is \texttt{article} with packages including geometry, babel, fontenc, setspace, enumitem, ulem, float, scrextend, xcolor, latexsym, amsmath, cancel, amssymb, amsthm; the wrapper is the lane's pinned \texttt{amsart} editorial wrapper at 10pt on A4, which restates the theorem-like environments the retained text needs and adds the article's own macro definitions that the retained text needs (\texttt{\textbackslash ex}), with extra wrapper packages (bm, bbm, dsfont, xspace, cancel, mathtools). The delivered numbering is the unit's own. Assets: no retained span contains a figure environment, a graphics-inclusion command, a PDF-page inclusion, a table or a TikZ picture; no image file is copied into this package, the package declares no asset dependency and no image file is redistributed with it. Licence: version arXiv:2401.10043v3 of this article is distributed under the Creative Commons Attribution 4.0 International licence, as recorded in the arXiv licence metadata for this exact version and shown on the abstract page https://arxiv.org/abs/2401.10043v3. A full rights notice is delivered at licenses/arxiv-2401.10043-CC-BY-4.0.txt. Characters: every retained excerpt is pure ASCII, so the delivered engine, XeLaTeX with its default Latin Modern text font, reports no missing character.
\begin{equation}\label{value_function}
    V(x) \coloneqq \inf\limits_{\substack{X(\cdot) \in \mathcal{A}(x) \\ \tau \in \mathcal{T}^X}} J(x, X(\cdot), \tau).
\end{equation}

\begin{equation}\label{value_function_constrained}
    V_\alpha(x, c) \coloneqq \inf\limits_{\substack{X(\cdot) \in \mathcal{A}(x) \\ \tau \in \mathcal{T}^X_\alpha}} J(x, X(\cdot), \tau, c).
\end{equation}

\begin{equation}\label{duality_gap}
    \begin{split}
        V_\alpha(x, c)
        &=
        \inf\limits_{\substack{X(\cdot) \in \mathcal{A}(x) \\ \tau \in \mathcal{T}^X}} 
        \left(
        \sup_{\lambda \ge 0}
        \Big(
            J(x, X(\cdot), \tau, c + \lambda) - \lambda \alpha 
        \Big)
        \right)
        \\&\ge
        \sup_{\lambda \ge 0}
        \left(
        \inf\limits_{\substack{X(\cdot) \in \mathcal{A}(x) \\ \tau \in \mathcal{T}^X}} 
        \Big(
            J(x, X(\cdot), \tau, c + \lambda) - \lambda \alpha 
        \Big)
        \right)
        \eqqcolon
        \Psi_\alpha(x, c),
    \end{split}
\end{equation}

\begin{Prop}\label{prop_duality_gap}
    Fix the parameters $x, c$ and $\alpha$.
    Suppose there exist a stopping time $\tau^\dag$ and a constant $\bar{\lambda} \ge 0$, such that the inequality $\ex[\tau^\dag] \le \alpha$ and the ``complementary slackness'' condition $\bar{\lambda}(\alpha - \ex[\tau^\dag]) = 0$ hold. 
    If
    \begin{equation}\label{prop_techn_equality}
        \inf\limits_{\substack{X(\cdot) \in \mathcal{A}(x) \\ \tau \in \mathcal{T}^X}}
            J\left(x, X(\cdot), \tau, c + \bar{\lambda}\right)
        =
        J\left(x, X^\dag(\cdot), \tau^\dag, c + \bar{\lambda}\right)
    \end{equation}
    also holds for some controlled process $X^\dag(\cdot) \in \mathcal{A}(x)$, i.e., if the pair $\left(X^\dag(\cdot), \tau^\dag\right)$ is optimal for the unconstrained problem of \eqref{value_function} with cost $c + \bar{\lambda}$ per unit of time, then $\left(X^\dag(\cdot), \tau^\dag\right)$ is optimal for the constrained problem \eqref{value_function_constrained}. Moreover, we have then the equality $\Psi_\alpha(x, c) = V_\alpha(x, c)$, meaning that there is no ``duality gap'' in \eqref{duality_gap}.
\end{Prop}

\begin{proof}
    First, we show that, under the assumptions of the proposition, the pair $(X^\dag(\cdot), \tau^\dag)$ is indeed optimal for the constrained problem \eqref{value_function_constrained}. Consider any admissible pair $(X(\cdot), \tau)$ for the constrained problem. Then the following chain of inequalities
    \begin{equation}
        \begin{split}
            J\left(x, X(\cdot), \tau, c\right)
            &\ge 
            J\left(x, X(\cdot), \tau, c + \bar{\lambda}\right) - \bar{\lambda} \alpha
            \\&\ge
            \inf\limits_{\substack{X(\cdot) \in \mathcal{A}(x) \\ \tau \in \mathcal{T}^X}} 
            J\left(x, X(\cdot), \tau, c + \bar{\lambda}\right) - \bar{\lambda} \alpha
            \\&=
            J\left(x, X^\dag(\cdot), \tau^\dag, c +\bar{\lambda}\right) - \bar{\lambda} \alpha
            \\&=
            J\left(x, X^\dag(\cdot), \tau^\dag, c\right)
        \end{split}
    \end{equation}
    proves the optimality of $(X^\dag(\cdot), \tau^\dag)$.
    Here, the first inequality follows from the fact that $\bar{\lambda} \ge 0$ and $\ex[\tau] \le \alpha$, the first equality is a consequence of the assumption \eqref{prop_techn_equality}, and the second equality is a consequence of the assumption $\bar{\lambda}(\alpha - \ex[\tau^\dag]) = 0$.

    The equality $\Psi_\alpha(x, c) = V_\alpha(x, c)$ is a consequence of the inequality \eqref{duality_gap} and of the reverse inequality
    \begin{equation}
    \begin{split}
        \Psi_\alpha(x, c) 
        &= 
        \sup_{\lambda \ge 0}
        \left(
        \inf\limits_{\substack{X(\cdot) \in \mathcal{A}(x) \\ \tau \in \mathcal{T}^X}} \Big(
            J(x, X(\cdot), \tau, c + \lambda) - \lambda \alpha 
        \Big)
        \right)
        \\&\ge
        \inf\limits_{\substack{X(\cdot) \in \mathcal{A}(x) \\ \tau \in \mathcal{T}^X}} \Big(
            J(x, X(\cdot), \tau, c + \bar{\lambda}) - \bar{\lambda} \alpha 
        \Big)
        \\&=
        J(x, X^\dag(\cdot), \tau^\dag, c + \bar{\lambda}) - \bar{\lambda} \alpha
        \\&=
        J(x, X^\dag(\cdot), \tau^\dag, c)
        \ge
        V_\alpha(x, c),
    \end{split}
    \end{equation}
    where the third equality follows again by the choice of $\bar{\lambda}$ and $\tau^\dag$.
\end{proof}

\end{document}
